\begin{figure}[htbp]
\centering
\begin{tikzpicture}[x=1cm,y=1cm,>=Stealth,font=\small]
\draw[green!65!black] (0,2.8)--(.7,3.4)--(.7,4.9)--(0,4.3)--cycle;
\node[red] at (.35,3.85) {图像};
\fill[yellow!85] (1.5,3.65) rectangle (3.2,4.25);\node[red] at (2.35,3.95) {CNN};
\draw[->,green!65!black] (.8,3.95)--(1.5,3.95);\draw[->,green!65!black] (3.2,3.95)--(4,3.95);
\draw[green!65!black] (4,2.8)--(4.8,3.5)--(4.8,5)--(4,4.2)--cycle;
\node[red] at (4.1,5.05) {特征图};\node[red] at (5.3,3.8) {CRF};
\draw[cyan!80!blue] (4.08,3.5)--(4.65,4.15) (4.25,3.48)--(4.25,4.2) (4.05,4.05)--(4.65,3.5) (4.16,3.72)--(4.58,3.74);
\foreach \xx/\yy in {4.22/3.7,4.55/3.72,4.52/4.2} \fill[red] (\xx,\yy) circle (2pt);
\node at (5.85,4.6) {节点};\draw[->,violet,dashed] (4.55,4.2)--(5.5,4.6);
\draw[cyan!80!blue] (6.1,3)--(6.1,3.7);\fill[red] (6.1,3) circle (2pt);\fill[red] (6.1,3.7) circle (2pt);
\node at (6,2.65) {CRF 图中的边（连接）};\draw[->,violet,dashed] (5.4,2.8)--(6.1,3.35);
\node at (-.3,2.75) {(a)};
\filldraw[fill=green!70!black,draw=magenta,dashed] (0,1) rectangle (1.9,2.4);
\draw[yellow] (.55,1.45)--(1,1.85);\fill[red] (.55,1.45) circle (3pt);\fill[red] (1,1.85) circle (3pt);
\filldraw[fill=green!70!black,draw=magenta,dashed] (2.3,.95) rectangle (4.2,2.35);
\draw[violet] (2.7,1.8)--(3.25,1.1)--(3.75,1.9);
\foreach \xx/\yy in {2.7/1.8,3.25/1.1,3.75/1.9} \fill[red] (\xx,\yy) circle (3pt);
\node at (-.3,.9) {(b)};
\draw[green!65!black] (4.55,1.35) rectangle (5.1,2);
\fill[red] (4.82,1.55) circle (2pt);\fill[red] (4.82,1.87) circle (2pt);
\draw[->,red,dashed] (4.55,3.72)--(4.82,2);
\fill[yellow!85] (5.5,.9) rectangle (6.8,2.25);\node[red,align=center] at (6.15,1.58) {全连接\\层};
\foreach \yy in {.9,1.3,1.65,2.25} \draw[->,green!65!black] (5.1,1.65)--(5.5,\yy);
\foreach \yy in {.95,1.6,2.1} \draw[->,magenta] (6.8,1.6)--(7.1,\yy);
\node[magenta,anchor=west] at (6.95,2.3) {$K\times K$};\node[magenta,anchor=west] at (7.15,1.6) {类别};
\node at (4.55,.8) {(c)};
\end{tikzpicture}
\origfig{11}{CRF 模型。（a）使用 CNN 从影像数据生成特征图；（b）空间范围框；（c）成对网络。}
\end{figure}
